В последние годы наблюдается активизация экспериментальных исследований,
направленных на повышение энергетической эффективности сегнерова колеса. География
этих работ охватывает как страны постсоветского пространства, так и государства Юго"=Восточной Азии, Латинской Америки и Европы, что свидетельствует о возобновившемся
интересе к реактивным турбинам для низконапорной гидроэнергетики. Одну из наиболее 
полных серий экспериментальных данных для диапазона напоров 2--5~м предоставляют 
исследования, выполненные на гидравлическом стенде в Ташкентском государственном 
техническом университете \cite{bozarov2023,bozarov2024a, bozarov2024b}. Конструкция испытанной установки представляет собой 
реактивно"=сопловую гидротурбину на основе сегнерова колеса с оптимизированной 
геометрией соплового аппарата. В ходе экспериментов фиксировались частота вращения, 
объёмный расход воды и развиваемая мощность при варьировании напора.

Во"=первых, экспериментальные значения КПД систематически превышают
теоретический предел $50\%$, установленный Эйлером для идеальной реактивной турбины.
Это превышение, достигающее 24~процентных пунктов при напоре 5~м, объясняется тем, 
что в реальной конструкции, в отличие от идеализированной модели, используются сопла
специальной профилированной формы, снижающие гидравлические потери, а также имеет
место частичное преобразование потенциальной энергии давления в подводящем тракте \cite{bozarov2024a}.

Во"=вторых, наблюдается монотонный рост КПД с увеличением напора, что указывает
на снижение доли относительных потерь (гидравлических и механических) при 
возрастании энергетического потенциала потока.

Принципиально иной уровень эффективности демонстрируют гибридные схемы,
сочетающие реактивную турбину с дополнительным активным рабочим органом. В такой
конфигурации вода, выходящая из сопел сегнерова колеса, попадает на лопатки встречно
вращающегося ротора, где отдаёт остаточную кинетическую энергию.

Экспериментальная проверка данной схемы \cite{bozarov2024b} при напоре 4,5~м показала КПД в
диапазоне $87{,}2$--$93{,}5\%$ в зависимости от режима нагружения. Сопоставление с данными
таблицы~\ref{tab:first} (КПД $72$--$74\%$ при сопоставимых напорах) позволяет утверждать, 
что переход к встречно"=роторной конструкции даёт прирост эффективности порядка 
$15$--20~процентных пунктов, что существенно превышает эффект от оптимизации геометрии 
сопел ($5$--$7\%$). Это открывает перспективы для использования сегнерова колеса не 
только как самостоятельной турбины для микроГЭС, но и как элемента более сложных 
гибридных энергетических установок, способных конкурировать с традиционными типами 
гидротурбин в низконапорном диапазоне.

Сопоставление теоретических расчётов с экспериментальными результатами \cite{bozarov2024a}
показывает, что расхождение составляет около $5{,}5\%$ при напоре 2~м и уменьшается 
с ростом напора. Это расхождение объясняется следующими факторами: 1) гидравлические потери в 
подводящем трубопроводе и соплах; 2) механические потери в подшипниках и уплотнениях;3) неравномерность 
поля скоростей на входе в сопла; 4) влияние вязкости жидкости на формирование струи. 

Одним из критических факторов, ограничивающих надёжность и долговечность
гидравлических турбин, является кавитация "--- явление образования и последующего
схлопывания паровых пузырьков в зонах пониженного давления. Для сегнерова колеса,
работающего на низких напорах (2--10~м), кавитационные риски считались минимальными,
однако современные исследования \cite{bozarov2024b} показывают, что при определённых режимах
кавитация может возникать и оказывать существенное влияние на эффективность и ресурс.

Гидравлические удары в момент замыкания сплошности вызывают эрозионное разрушение 
материала гребных винтов, роторов насосов и~т.~п., а также специфическую коррозию 
металлов, из которых сделаны эти детали, из"=за того, что пассивирующая оксидная плёнка 
непрерывно удаляется. При невозможности избежать кавитации на практике её стараются 
усилить и перевести в режим «суперкавитации» и струйного обтекания.

Если пузырёк содержит мало газа, он схлопывается полностью в первом же периоде. 
Сокращение пузырька происходит с большой скоростью и сопровождается звуковым импульсом, 
который тем сильнее, чем меньше газа было в пузырьке. При высокой интенсивности 
кавитации образуется и схлопывается множество пузырьков, которые создают сильный шум 
сплошного спектра в диапазоне от сотен герц до сотен килогерц. Обтекаемое жидкостью 
тело оказывается окружено чётко ограниченной кавитационной зоной, заполненной 
движущимися пузырьками.

Кавитация возникает в тех областях проточной части, где статическое давление
падает ниже давления насыщенных паров жидкости. Для воды при температуре 
$20\,^{\circ}\text{C}$ давление насыщенных паров составляет $\sim 2{,}34$~кПа (0,0234~атм). 
В сегнеровом колесе зонами потенциального кавитационного риска являются:

\begin{enumerate}
    \renewcommand{\labelenumi}{\arabic{enumi})}
    \item входные кромки сопел, где скорость потока максимальна, а статическое давление минимально;
    \item выходные сечения сопел, где происходит резкое расширение струи, сопровождающееся образованием зон пониженного давления;
    \item поверхности вращающегося ротора в зоне высоких окружных скоростей (при $u > 10$--$12$~м/с).
\end{enumerate}

Также количественная оценка условий, при которых начинается кавитация, может быть 
выполнена на основе энергетического баланса потока. Кавитация возникает, если 
выполняется условие:

\begin{equation*}
    P + \frac{\rho V^2}{2} > Z,
\end{equation*}
где \texttt{P} "--- гидростатическое давление в потоке, \texttt{$\rho$} "= плотность жидкости, \texttt{V} "= скорость потока, \texttt{Z} "=
объёмная прочность жидкости, которую в первом приближении можно принимать равной
давлению насыщенных паров $P_{\text{np}}$. Из этого условия критическая скорость, при которой
начинается кавитация, определяется как:
\begin{equation*}
    V_{\text{кр}} = \sqrt{\frac{P_{\text{np}} - P}{\rho}}
\end{equation*}

Для воды при атмосферном давлении ($P = 101{,}3$~кПа) и температуре $20\,^{\circ}\text{C}$ 
($P_{\text{нп}} = 2{,}34$~кПа) критическая скорость составляет около 14~м/с. Однако, 
как показывают экспериментальные данные, для реальных конструкций с учётом геометрии 
проточных каналов это значение может снижаться до 7--8~м/с. Это объясняется влиянием 
формы внутренних поверхностей на условия зарождения кавитационных пузырьков, что 
необходимо учитывать при проектировании соплового аппарата.

\lstinputlisting[caption={Код программы расчёта критической скорости кавитации}, label={lst:cod}]{cod.cpp}

В реальных условиях, с учётом гидравлических потерь, действительная скорость
вытекания меньше теоретической и определяется с использованием коэффициента скорости 
$\varphi$:
\begin{equation*}
    v_{\text{д}} = \varphi \sqrt{2gH\nu}
\end{equation*}

Коэффициент скорости зависит от геометрии сопла, шероховатости внутренней
поверхности и вязкости жидкости. Экспериментальные данные показывают, что для 
хорошо спрофилированных сопел $\varphi$ находится в диапазоне $0{,}95$--$0{,}98$ \cite{bozarov2024a}.
